\documentclass[12pt,a4paper]{article}

% =========================================================================
% FORMATAÇÃO ABNT (NBR 14724 / NBR 6024 / NBR 6027)
% =========================================================================

\usepackage[
  a4paper,
  left=3cm,
  top=3cm,
  right=2cm,
  bottom=2cm,
  headheight=15pt,
  headsep=1cm
]{geometry}

\usepackage[T1]{fontenc}
\usepackage[utf8]{inputenc}
\usepackage[brazil]{babel}
\usepackage{amsmath}
\usepackage{newtxtext,newtxmath}
\usepackage{microtype}
\usepackage[htt]{hyphenat}

\usepackage{setspace}
\onehalfspacing
\usepackage{indentfirst}
\setlength{\parindent}{1.25cm}
\setlength{\parskip}{0pt}
\setlength{\emergencystretch}{2em}

\usepackage{array}
\usepackage{tabularx}
\usepackage{booktabs}
\usepackage{etoolbox}

\usepackage{graphicx}
\graphicspath{{../assets/}{assets/}}
\usepackage{float}

\usepackage{fancyhdr}
\pagestyle{empty}
\fancyhf{}
\fancyhead[R]{\small\thepage}
\renewcommand{\headrulewidth}{0pt}

\usepackage{tocloft}
\renewcommand{\cftsecleader}{\cftdotfill{\cftdotsep}}
\renewcommand{\cftsecfont}{\bfseries}
\renewcommand{\cftsecpagefont}{\bfseries}
\setlength{\cftbeforesecskip}{0.4em}

\usepackage{titlesec}
\titleformat{\section}
  {\normalfont\normalsize\bfseries\MakeUppercase}{\thesection}{1em}{}
\titleformat{\subsection}
  {\normalfont\normalsize\bfseries}{\thesubsection}{1em}{}
\titleformat{\subsubsection}
  {\normalfont\normalsize\bfseries}{\thesubsubsection}{1em}{}
\titleformat{\paragraph}
  {\normalfont\normalsize\itshape}{\theparagraph}{1em}{}

\titlespacing*{\section}{0pt}{2ex plus 0.5ex minus 0.2ex}{1.5ex plus 0.5ex minus 0.2ex}
\titlespacing*{\subsection}{0pt}{1.8ex plus 0.5ex minus 0.2ex}{1.2ex plus 0.5ex minus 0.2ex}
\titlespacing*{\subsubsection}{0pt}{1.5ex plus 0.5ex minus 0.2ex}{1.0ex plus 0.5ex minus 0.2ex}

\usepackage{caption}
\captionsetup{
  font=small,
  labelfont=bf,
  justification=centering,
  singlelinecheck=false,
  skip=6pt
}
\captionsetup[figure]{name=Figura, labelsep=endash}
\captionsetup[table]{name=Quadro, labelsep=endash}

\newcommand{\fonte}[1]{%
  \par\vspace{2pt}%
  {\footnotesize\raggedright\textbf{Fonte:} #1\par}%
}

\AtBeginEnvironment{figure}{\singlespace}
\AtBeginEnvironment{table}{\singlespace}

\usepackage{xcolor}
\usepackage{listings}
\ifdefined\pdftexversion
  \usepackage{listingsutf8}
\fi
\definecolor{codgray}{gray}{0.45}
\lstset{
  basicstyle=\ttfamily\footnotesize,
  breaklines=true,
  columns=fullflexible,
  keepspaces=true,
  showstringspaces=false,
  frame=single,
  framexleftmargin=4pt,
  tabsize=2,
  numbers=left,
  numberstyle=\tiny\color{codgray},
  numbersep=8pt,
  xleftmargin=20pt,
  captionpos=b
}
\ifdefined\pdftexversion
  \lstset{inputencoding=utf8/latin1, extendedchars=true}
\fi

\providecommand{\tightlist}{%
  \setlength{\itemsep}{0pt}\setlength{\parskip}{0pt}}

\usepackage[hidelinks]{hyperref}
\usepackage{bookmark}
\IfFileExists{xurl.sty}{\usepackage{xurl}}{}
\urlstyle{same}

\begin{document}

% =========================================================================
% ELEMENTOS PRÉ-TEXTUAIS
% =========================================================================

% --- CAPA ---
\begin{titlepage}
  \thispagestyle{empty}
  \begin{center}
    \textbf{\large UNIVERSIDADE FEDERAL DE GOIÁS}\\[0.2cm]
    \textbf{\large INSTITUTO DE INFORMÁTICA}\\[0.2cm]
    \textbf{\large BACHARELADO EM ENGENHARIA DE SOFTWARE}\\[0.2cm]
    \textbf{\large DISCIPLINA: MODELAGEM DE SOFTWARE}

    \vspace{3.2cm}
    {\large
    Gabriel Henrique Lourenço Dias\par}

    \vspace{3.5cm}
    \textbf{\Large SITE DE LOJA DE ELETRÔNICOS}\\[0.4cm]
    \textbf{\large Especificação de Requisitos --- APS 2}\par

    \vfill
    \textbf{GOIÂNIA}\\[0.2cm]
    \textbf{2026}
  \end{center}
\end{titlepage}

% --- FOLHA DE ROSTO ---
\thispagestyle{empty}
\begin{center}
  {\large
  Gabriel Henrique Lourenço Dias\par}

  \vspace{4.0cm}
  \textbf{\Large SITE DE LOJA DE ELETRÔNICOS}\\[0.4cm]
  \textbf{\large Especificação de Requisitos --- APS 2}\par
\end{center}

\vspace{2.5cm}
\begin{flushright}
  \begin{minipage}{0.52\textwidth}
    \begin{singlespace}
      {\small Trabalho acadêmico apresentado à disciplina de Modelagem
      de Software do curso de Bacharelado em Engenharia de Software da
      Universidade Federal de Goiás (UFG), como requisito parcial de
      avaliação referente à Atividade Prática Supervisionada (APS
      2).\\[0.4cm]
      \textbf{Professor Orientador:} Prof. Me. Diogo Oliveira Santos}
    \end{singlespace}
  \end{minipage}
\end{flushright}

\vfill
\begin{center}
  \textbf{GOIÂNIA}\\[0.2cm]
  \textbf{2026}
\end{center}
\newpage

% --- SUMÁRIO ---
\thispagestyle{empty}
\renewcommand{\contentsname}{SUMÁRIO}
\pdfbookmark[0]{\contentsname}{toc}
\tableofcontents
\newpage

% =========================================================================
% ELEMENTOS TEXTUAIS
% =========================================================================

\pagestyle{fancy}
\thispagestyle{fancy}

\section{Visão Geral do Sistema}\label{sec:visao-geral}

\subsection{Objetivo}

Disponibilizar um site de loja de eletrônicos em que o cliente navega
por um catálogo (com filtros por categoria, marca e preço), monta um
carrinho e finaliza o pedido, com pagamento via gateway e cálculo de
frete por serviço externo. Um módulo administrativo permite cadastrar
produtos, gerenciar o estoque e acompanhar o status dos pedidos.

\subsection{Problema que Resolve}

A venda de eletrônicos é limitada ao atendimento presencial ou por
canais informais, o controle de estoque é manual gerando venda de itens
indisponíveis, e não há acompanhamento do pedido por parte do cliente.
O sistema centraliza o catálogo, valida disponibilidade em tempo real e
oferece rastreamento de pedidos com pagamento seguro via gateway
externo.

\subsection{Benefícios Esperados}

\begin{itemize}
\tightlist
\item Venda 24/7 sem depender de atendente
\item Descoberta rápida de produtos por filtros
\item Estoque consistente com reserva e baixa automáticas
\item Pagamento seguro (dados tokenizados)
\item Frete transparente exibido antes da confirmação
\item Rastreabilidade do status do pedido
\end{itemize}

\section{Atores}\label{sec:atores}

O sistema conta com os seguintes atores:

\begin{itemize}
\tightlist
\item \textbf{Visitante:} navega no catálogo, filtra, busca, vê detalhes,
  monta carrinho, calcula frete e se cadastra (sem autenticação).
\item \textbf{Cliente:} visitante autenticado que finaliza pedidos, paga,
  acompanha e cancela pedidos (seus próprios).
\item \textbf{Administrador:} funcional da loja que gerencia produtos,
  categorias, marcas, estoque e status dos pedidos.
\item \textbf{Gateway de Pagamento} (externo): autoriza, captura e estorna
  pagamentos.
\item \textbf{Serviço de Frete} (externo): calcula valor e prazo de entrega.
\item \textbf{Serviço de E-mail} (externo): entrega e-mails transacionais.
\end{itemize}

\section{Requisitos Funcionais}\label{sec:rf}

\subsection{Catálogo}

\begin{itemize}
\tightlist
\item RF-01: Listar produtos ativos com paginação
\item RF-02: Filtrar por categoria, marca e faixa de preço
\item RF-03: Buscar por texto
\item RF-04: Ordenar por preço, nome ou relevância
\item RF-05: Exibir detalhes completos do produto
\end{itemize}

\subsection{Conta e Autenticação}

\begin{itemize}
\tightlist
\item RF-06: Cadastrar cliente (nome, CPF, e-mail, telefone, senha)
\item RF-07: Autenticar Cliente e Administrador (e-mail + senha)
\item RF-08: Recuperar senha (link por e-mail)
\item RF-09: Gerenciar endereços de entrega
\end{itemize}

\subsection{Carrinho}

\begin{itemize}
\tightlist
\item RF-10: Adicionar ao carrinho com validação de disponibilidade
\item RF-11: Alterar quantidade ou remover item
\item RF-12: Visualizar carrinho com subtotal
\item RF-13: Manter carrinho ao autenticar (visitante → cliente)
\end{itemize}

\subsection{Frete}

\begin{itemize}
\tightlist
\item RF-14: Calcular opções de frete por CEP de destino
\end{itemize}

\subsection{Pedido e Pagamento}

\begin{itemize}
\tightlist
\item RF-15: Finalizar pedido (checkout)
\item RF-16: Registrar pedido com número único
\item RF-17: Processar pagamento no gateway
\item RF-18: Receber confirmação assíncrona (Pix, boleto)
\item RF-19: Acompanhar pedido (status e histórico)
\item RF-20: Cancelar pedido (com estorno se pago)
\item RF-21: Estornar pagamento
\item RF-22: Notificar por e-mail (eventos do pedido)
\end{itemize}

\subsection{Estoque}

\begin{itemize}
\tightlist
\item RF-23: Verificar disponibilidade
\item RF-24: Reservar e baixar estoque
\item RF-25: Gerenciar estoque (entrada/saída/ajuste)
\item RF-26: Alertar estoque baixo
\end{itemize}

\subsection{Módulo Administrativo}

\begin{itemize}
\tightlist
\item RF-27: Cadastrar produto (SKU, nome, descrição, preço, estoque)
\item RF-28: Editar ou inativar produto
\item RF-29: Gerenciar categorias e marcas
\item RF-30: Consultar pedidos
\item RF-31: Atualizar status do pedido
\item RF-32: Cancelar pedido (com motivo)
\end{itemize}

\section{Requisitos Não Funcionais}\label{sec:rnf}

\subsection{Segurança}

\begin{itemize}
\tightlist
\item RNF-01: Senhas com hash (bcrypt/argon2)
\item RNF-02: HTTPS/TLS 1.2+
\item RNF-03: Controle de acesso por perfil (RBAC)
\item RNF-04: Bloqueio temporário após 5 falhas de login
\item RNF-05: Sem armazenamento de número de cartão (tokenização)
\item RNF-06: Proteção contra ataques OWASP Top 10
\end{itemize}

\subsection{Usabilidade}

\begin{itemize}
\tightlist
\item RNF-07: Interface responsiva (320px a 1920px)
\item RNF-08: Checkout em máximo 4 etapas
\item RNF-09: Mensagens de erro claras, acessibilidade WCAG AA
\end{itemize}

\subsection{Desempenho}

\begin{itemize}
\tightlist
\item RNF-10: Busca ≤ 500ms (P95) em 200 usuários simultâneos
\item RNF-11: Carregamento ≤ 3s (P90)
\item RNF-12: Serviços externos assíncronos com timeout e fallback
\end{itemize}

\subsection{Disponibilidade e Escalabilidade}

\begin{itemize}
\tightlist
\item RNF-13: 99,5\% de disponibilidade mensal
\item RNF-14: Backup diário com retenção de 30 dias
\item RNF-15: Falha em serviço externo degrada apenas essa função
\item RNF-16: Backend stateless (réplicas horizontais)
\item RNF-17: Cache de catálogo (TTL curto)
\item RNF-18: Arquitetura modular com interfaces bem definidas
\end{itemize}

\subsection{Privacidade e Auditoria}

\begin{itemize}
\tightlist
\item RNF-19: Dados pessoais conforme LGPD
\item RNF-20: Log de auditoria das ações administrativas
\item RNF-21: Sem dados pessoais em URLs ou logs
\end{itemize}

\section{Regras de Negócio}\label{sec:rn}

\subsection{Catálogo}

\begin{itemize}
\tightlist
\item RN-01: Apenas produtos \texttt{ATIVO} aparecem no catálogo
\item RN-02: Produto sem estoque marca-se indisponível (não pode adicionar ao carrinho)
\item RN-03: Filtros combinados usam E entre critérios distintos
\item RN-32: SKU único; preço > 0; exatamente 1 categoria e 1 marca
\item RN-33: Exige ao menos 1 imagem, descrição, peso e dimensões
\item RN-34: Produto com pedidos não é excluído, apenas inativado
\item RN-35: Categoria/marca só é inativada se nenhum produto ativo a usar
\end{itemize}

\subsection{Acesso e Conta}

\begin{itemize}
\tightlist
\item RN-04: CPF válido e único
\item RN-05: E-mail válido e único
\item RN-06: Senha ≥ 8 caracteres, ≥1 letra, ≥1 número
\item RN-07: Perfil derivado do cadastro (exclusivo)
\item RN-08: Bloqueio de 15min após 5 falhas de login
\item RN-09: Link de senha expira em 30min, uso único
\item RN-10: Pedido exige endereço completo com CEP válido
\end{itemize}

\subsection{Carrinho e Frete}

\begin{itemize}
\tightlist
\item RN-11: Quantidade ≥1 e ≤ estoque disponível
\item RN-12: Disponível = em estoque − reservado
\item RN-13: Total = soma(preço × qtd) + frete, congelado na criação do pedido
\item RN-14: Itens do carrinho visitante somam-se ao do cliente ao autenticar
\item RN-15: Frete calculado por CEP destino + peso/dimensões dos itens
\item RN-16: Sem frete disponível, checkout não prossegue
\end{itemize}

\subsection{Pedido}

\begin{itemize}
\tightlist
\item RN-17: Somente cliente autenticado finaliza pedido
\item RN-18: Frete escolhido confirmado; mudança invalida o cálculo
\item RN-19: Pedido nasce com status \texttt{AGUARDANDO\_PAGAMENTO}
\item RN-20: Formas aceitas: cartão (imediato), Pix e boleto (pendentes)
\item RN-21: Recusa não cancela pedido; cliente tenta outra forma
\item RN-22: Pagamento aprovado → status \texttt{PAGO}
\item RN-23: Cliente vê seus pedidos; admin vê todos
\item RN-24: Cancelamento permitido enquanto não \texttt{ENVIADO}
\item RN-25: Pedido pago e cancelado → estorno total
\item RN-26: Falha de e-mail não desfaz operação; reprocessa-se
\end{itemize}

\subsection{Estoque}

\begin{itemize}
\tightlist
\item RN-27: Estoque reservado ao criar pedido, baixado ao confirmar pagamento
\item RN-28: Reserva devolvida se pedido cancelado ou expirar
\item RN-29: Estoque nunca fica negativo
\item RN-30: Ajuste manual registra motivo, autor, data
\item RN-31: Produtos abaixo do mínimo geram alerta
\end{itemize}

\subsection{Administração}

\begin{itemize}
\tightlist
\item RN-36: Mudanças registram autor e data em auditoria; envio exige rastreio
\end{itemize}

\section{Casos de Uso}\label{sec:casos-uso}

O sistema contém 23 casos de uso organizados em:

\begin{itemize}
\tightlist
\item \textbf{Catálogo:} navegar, filtrar, buscar, ver detalhes
\item \textbf{Conta:} cadastro, autenticação, recuperação de senha, endereços
\item \textbf{Carrinho:} adicionar, alterar, visualizar, manter sessão
\item \textbf{Frete:} calcular opções
\item \textbf{Pedido:} finalizar, acompanhar, cancelar
\item \textbf{Pagamento:} processar (cartão, Pix, boleto), estornar, confirmação assíncrona
\item \textbf{Notificação:} enviar e-mails de eventos
\item \textbf{Estoque:} reservar, baixar, gerenciar, alertar
\item \textbf{Administrativo:} produtos, categorias, marcas, pedidos, status
\end{itemize}

Relacionamentos propostos: generalização de Visitante/Cliente e de
Pagamento (especializado em cartão, Pix, boleto); inclusão obrigatória
em fluxos críticos (autenticação, frete, pagamento); extensão em cenários
opcionais (filtros, busca, estorno condicional).

\section{Arquitetura e Componentes}\label{sec:arquitetura}

\subsection{Estilo}

Monólito modular com front-ends separados e serviços externos via
adaptadores. Cada módulo de negócio expõe interface e consome outras
apenas por interface, permitindo trocar fornecedores externos.

\subsection{Pacotes}

\begin{itemize}
\tightlist
\item \textbf{Frontend:} interface web da loja; painel administrativo
\item \textbf{Backend -- Núcleo:} API, usuários, catálogo, busca, carrinho,
  pedidos, estoque, notificações
\item \textbf{Backend -- Integrações:} adaptadores de pagamento, frete, e-mail
\item \textbf{Dados:} banco de dados relacional
\item \textbf{Serviços Externos:} gateway de pagamento, frete, e-mail
\end{itemize}

\subsection{Componentes Principais}

\begin{itemize}
\tightlist
\item \textbf{API de Aplicação:} ponto de entrada HTTPS; roteia, valida
  sessão e orquestra módulos
\item \textbf{Módulo de Catálogo:} produtos, categorias, marcas
\item \textbf{Módulo de Pedidos:} checkout, status, orquestração
\item \textbf{Módulo de Estoque:} disponibilidade, reserva, ajustes
\item \textbf{Módulo de Notificações:} eventos, fila, reprocessamento
\item \textbf{Adaptadores:} pagamento, frete, e-mail (interface com externos)
\end{itemize}

\section{Premissas e Decisões}\label{sec:premissas}

\begin{itemize}
\tightlist
\item Web app responsivo (não nativo)
\item Cadastro/autenticação acrescentados (não constam na descrição)
\item Um único perfil administrativo
\item Formas de pagamento: cartão, Pix, boleto
\item Visitante pode navegar sem autenticar; autenticação só no checkout
\item Notificações apenas por e-mail
\item Cancelamento permitido até separação; envio não cancela
\item Pedidos em espera por pagamento expiram (configurável, ex.: 24h)
\item Arquitetura em monólito modular com adaptadores
\item Sem cupons, promoções, devolução, nota fiscal, marketplace nesta versão
\end{itemize}

\section{Rastreabilidade}\label{sec:rastreabilidade}

A especificação fornece base para os dois diagramas exigidos:

\begin{itemize}
\tightlist
\item \textbf{Diagrama de Casos de Uso:} seção \ref{sec:atores} (atores e
  generalização) e seção \ref{sec:casos-uso} (23 UCs com relacionamentos
  de inclusão e extensão).
\item \textbf{Diagrama de Componentes:} seção \ref{sec:arquitetura}
  (5 pacotes, 13 componentes, interfaces fornecidas/requeridas, dependências,
  portas).
\end{itemize}

Todos os requisitos funcionais (RF), não funcionais (RNF) e regras de
negócio (RN) estão rastreados na especificação completa em
\texttt{aps2/docs/especificacao-requisitos.md}.

\end{document}
